\scanpage{114}

% manuscript lines / position: 1–3
\uline{NB} la condition sur $N_2$ implique celle sur $T_2$, en passant aux formes bil. ass., i.e. en l'appliquant pour $x$ \uline{et} pour $1+x$ $(x\in A_1)$.

% manuscript lines / position: 4
\textcircled{$3^\circ$} On voit la condition

% manuscript lines / position: 5
\[
 \boxed{\sigma_1\sigma_2\sigma_3=1}
\]

% manuscript lines / position: 4–9, left, oblique
% slanted
\margin{Cette condition est automatiquement vérifiée sur la sous-algèbre de $A$ engendrée par les $A_i$, car elle est vraie sur \unsure{chaque|les} $A_i$ (lequel est stable par les $\sigma_i$).}

% manuscript lines / position: 6–8
implique que les involutions $\sigma_i$ commutent entre elles, car le produit de 2 d'entre elles est une involution

% manuscript lines / position: 8–10
\[
 \boxed{\begin{aligned}\sigma_1\sigma_2&=\sigma_3\\
 \sigma_2\sigma_3&=\sigma_1\\
 \unsure{\cdot}\sigma_3\sigma_1&=\sigma_2\end{aligned}}
\]

% manuscript lines / position: 11–13
(car pour deux inv. $\sigma$ et $\tau$,

% manuscript lines / position: 12–13
\[
 \begin{aligned}
 \sigma\tau=\tau\sigma &\iff (\sigma\tau)^2=\deleted{\unsure{\sigma}}\;1\quad\text{i.e.}\\
 &\hphantom{\iff}\sigma\tau\sigma\tau=1
 \end{aligned}\quad\text{)}
\]

% manuscript lines / position: 14–15
En utilisant $1^\circ)$, \deleted{\unsure{ceci}} \inserted{la condition $\sigma_1\sigma_2\sigma_3=1$} s'explicite par la formule

% manuscript lines / position: 16
\[
 2x=T_1(x)+T_2(x)+T_3(x)-\underline{t_2T_{2}(x)}\unsure{))},
\]

% manuscript lines / position: 17–18
en apparence asymétrique, \deleted{mais} elle est équivalente à $\sigma_2\sigma_3\sigma_1=1$ ou $\sigma_3\sigma_1\sigma_2=1$,

% manuscript lines / position: 19–20
\[
 \begin{aligned}
 \text{donc}\qquad 2x&=T_1(x)+T_2(x)+T_3(x)-t_3T_3(x)\\
 2x&=\qquad -\;\;-\;\;-\;\;\text{——}\quad -t_1T_1(x),
 \end{aligned}
\]

% manuscript lines / position: 21
donc elle \uline{implique}\textsuperscript{*} la condition

% manuscript lines / position: 21–24, left, oblique
% slanted
\margin{du moins si $k\longrightarrow A$ injectif --- mais on la \unsure{supposera} vérifiée en tout cas\ldots}

% manuscript lines / position: 22
\[
 \boxed{t_1\mathrm{Tr}_1(x)=t_2\mathrm{Tr}_2(x)=t_3\mathrm{Tr}_3(x)}
\]

% manuscript lines / position: 23–25
(qui ressemble à $6^\circ$ ci-dessus). \unsure{D'on|D'où|Si on} \unsure{désigner|désignant|désigne} par $T_{A/k}(x)$ la valeur commune aux trois mbres, on trouve donc la condition

% manuscript lines / position: 26
\[
 \boxed{2x=T_1(x)+T_2(x)+T_3(x)-T_{A/k}(x).}
\]

% manuscript lines / position: bottom left
% archive sheet number: 113
